% source: d666a7a4f8b4ae6d33990de81fc5d0e4fb6ddafc
\chapter{最短向量问题（SVP）}
\label{ch:svp}

\section{什么是最短向量？}

回忆一下，给定一组基
\[
B=(b_1,b_2,\dots,b_n),
\]
格由其所有整系数线性组合构成：
\[
L(B)=\left\{\sum_{i=1}^n z_i b_i : z_i\in\mathbb{Z}\right\}.
\]

最短向量问题（shortest vector problem, SVP）要求在所有非零格点中找出欧几里得范数最小的那一个。形式化地，
\[
\lambda_1(L)=\min_{v\in L\setminus\{0\}} \|v\|.
\]
量 $\lambda_1(L)$ 称为第一逐次极小（first successive minimum）。在文献中，这个量在格密码学里反复出现。

换言之，如果我们设想基向量的所有整数组合，那么最短向量就是离原点最近的非零格点。

\begin{figure}[H]
\centering
\begin{tikzpicture}[scale=0.55]
  \draw[pqcgray!40, thin] (-5.5,0) -- (5.5,0);
  \draw[pqcgray!40, thin] (0,-6.5) -- (0,6.5);
  \foreach \i in {-2,...,2} {
    \foreach \j in {-2,...,2} {
      \pgfmathsetmacro{\x}{2*\i+1*\j}
      \pgfmathsetmacro{\y}{1*\i+3*\j}
      \node[latticept] at (\x,\y) {};
    }
  }
  \node[latticeptbold] at (0,0) {};
  \draw[shortvec] (0,0) -- (1,-2);
  \node[pqcorange, font=\scriptsize] at (1.55,-1.9) {$\lambda_1$};
\end{tikzpicture}
\caption{与第~\ref{ch:lattice}~章中相同的格 $L(b_1,b_2)$，其中 $b_1=(2,1)$，$b_2=(1,3)$。在图中所有点里，$(1,-2)=b_1-b_2$ 是一个最短非零向量，$\lambda_1(L)=\sqrt{5}$。注意，它是基向量的一个\emph{组合}，而不是某个基向量本身。}
\label{fig:svp}
\end{figure}

\section{为什么二维 SVP 很容易？}

在低维情形下，几何结构仍然是可见的。在二维中，人们往往可以看着格的图形用肉眼找出最短向量。在三维中仍可进行几何推理，但图形已不那么直观。到了四维及以上，人的直觉变得很弱。等到 256 维时，已经没有任何有意义的方式能直接将其几何结构可视化。

这就是格问题在高维中变得计算困难的原因。搜索空间呈指数增长，蛮力方法毫无希望。

\section{为什么蛮力搜索不可行}

假设我们尝试在维数 $100$ 的格中进行搜索。若每个系数取自 $-5,\dots,5$，则每个坐标有 $11$ 种可能。候选向量的总数为
\[
11^{100},
\]
约为
\[
1.37\times 10^{104}.
\]
这远远超过可观测宇宙中的原子数。因此穷举搜索是不可行的。

确切地说，蛮力搜索就是第~\ref{ch:lattice}~章中的算法~\ref{alg:lattice-enum}，外加一个记录迄今所见最小范数的步骤：

\begin{algorithm}[H]
\caption{蛮力最短向量搜索}
\label{alg:svp-bruteforce}
\begin{algorithmic}[1]
\Require 基 $b_1,\dots,b_n$；系数界 $N$
\Ensure $\{\sum_i z_i b_i : |z_i|\le N,\ z\ne 0\}$ 中的一个最短向量
\State $v^\star \gets \textsc{null}$, \quad $m \gets \infty$
\For{每个 $(z_1,\dots,z_n) \in \{-N,\dots,N\}^n \setminus \{0\}$}
    \State $v \gets \sum_{i=1}^{n} z_i\, b_i$
    \If{$\|v\| < m$}
        \State $v^\star \gets v$, \quad $m \gets \|v\|$
    \EndIf
\EndFor
\State \Return $v^\star$
\end{algorithmic}
\end{algorithm}

循环体的开销很小，但循环本身要访问 $(2N+1)^n$ 个候选。当 $n=100$、$N=5$ 时，这正是上面的 $11^{100}$；使精确 SVP 变得困难的，是这种随 $n$ 的指数爆炸，而不是内层计算的任何低效。

\section{初学者的常见陷阱}

许多初学者以为最短向量必定是某个基向量本身。这并不正确。最短向量是一个格向量，但它不必等于任何单个基向量。例如，最短向量可以是基向量的组合，如
\[
 b_2-b_1,
\]
它可能比任一基向量单独都更短。

这意味着求解 SVP 需要考察整个格结构，而不仅仅是恰好给定的那组基。

\section{Sage 实验}

我们可以构造一个小格，并让 Sage 求出它的最短向量。

\begin{lstlisting}[language=Python]
from sage.modules.free_module_integer import IntegerLattice

B = matrix(ZZ, [[2, 1], [1, 3]])
L = IntegerLattice(B)
print(L.shortest_vector())
print(L.shortest_vector().norm())
\end{lstlisting}

输出给出一个最短格向量及其欧几里得范数。这就是最基本的 SVP 实例。

\section{Minkowski 定理}

Minkowski 定理是格理论的基础性结果之一。它指出，$n$ 维格必定包含一个非零向量，其范数至多为
\[
\sqrt{n}\,\det(L)^{1/n}.
\]

其中的重要直觉很简单：如果格是稠密的，那么最短向量不可能太长；如果格是稀疏的，那么最近的格点就离得更远。Minkowski 定理保证了存在一个不太长的向量，尽管它并没有告诉我们如何高效地找到它。

\section{精确 SVP 与近似 SVP}

精确 SVP 与近似 SVP 之间有一个重要区别。

\begin{itemize}
  \item 精确 SVP 要求找到真正的最短向量。
  \item 近似 SVP 只要求找到一个长度在最优值的 $\gamma$ 倍以内的向量。
\end{itemize}

按照文献中的记法，$\gamma$-SVP 求解器返回的向量长度至多为
\[
\gamma\,\lambda_1(L).
\]

这一区别很重要，因为格密码分析中的实用算法通常并不求解精确 SVP，而是求解近似变体，后者要容易处理得多。

\section{为什么 LLL 不能求解精确 SVP}

LLL 算法极其有用，但它的目标并不是找出精确的最短向量。它所做的是把一组基变换为约化程度更高的基，往往能产生一个比原始向量短得多的向量。然而，LLL 返回的向量通常并不保证是真正的最短向量。在许多情况下，它只是一个良好的近似。

\section{为什么 BKZ 更强}

BKZ 通过处理格的局部分块来改进 LLL。它在一个滑动窗口内反复求解小维数的 SVP 实例，然后利用所得的短向量来改进全局基。这使 BKZ 比 LLL 更强，但也大幅增加了计算开销。在实践中，BKZ 能找到非常短的向量，但在高维中仍然不能保证求解精确 SVP。

\section{为什么 SVP 是困难的}

就目前而言，已知求解精确 SVP 的最佳经典算法在维数上仍是指数级的，而且没有任何已知的量子算法能提供 Shor 算法对因数分解所带来的那种巨大加速。因此，基于格的密码仍是实现后量子安全的一条实用且有理论依据的途径。

\section{与 LWE 的联系}

与带误差学习（LWE）的联系是核心所在。LWE 问题可以看作格上的一个带噪声的线性方程组。如果能够高效求解相关的格问题（包括在适当设定下的 SVP），那么许多针对 LWE 的攻击将变得强大得多。Regev 归约的方向则相反：最坏情形的格问题可归约到平均情形的 LWE，因此一个能在随机实例上攻破 LWE 的高效算法将能在\emph{每一个}格上求解近似格问题。第~\ref{ch:lwe}~章将精确陈述这一定理，并讨论它在实际部署参数下能保证什么、不能保证什么。

\section{要点}

\begin{itemize}
    \item SVP 要求的是最短的非零格向量，而不仅仅是最短的基向量。
    \item SVP 的困难性随维数迅速增长。
    \item LLL 与 BKZ 求解的是 SVP 的近似版本，一般而言并不求解精确问题。
    \item SVP 与 LWE 之间的联系是人们相信基于格的密码安全的核心原因之一。
\end{itemize}

\printbibliography[heading=chapterrefs]
